% ==================================================
% MECÁNICA DE DISTRIBUCIONES — archivo principal
% Este archivo solo "monta" el libro; el contenido
% vive en capitulos/, apendices/ y portada/.
% ==================================================
% Para IMPRESIÓN con capítulos siempre en página impar: cambia openany por openright
\documentclass[11pt,a4paper,openany]{book}

% ==================================================
% PREÁMBULO DEL LIBRO
% Aquí van paquetes, colores, comandos y entornos.
% Normalmente NO hace falta tocar este archivo para escribir.
% ==================================================

% Idioma y codificación
\usepackage[utf8]{inputenc}
\usepackage[T1]{fontenc}
\usepackage[spanish,es-nodecimaldot]{babel}

% Tipografía
\usepackage{newtxtext}
\usepackage{newtxmath}
\usepackage{microtype}

% Página
\usepackage[
    a4paper,
    inner=2.8cm,
    outer=2.3cm,
    top=2.5cm,
    bottom=2.5cm,
    headheight=15pt
]{geometry}

% Matemáticas
\usepackage{amsmath,amssymb,amsthm,mathtools}
\usepackage{bm}

% Gráficos y tablas
\usepackage{graphicx}
\usepackage{booktabs}
\usepackage{array}

% Color y cajas
\usepackage[dvipsnames,svgnames]{xcolor}
\usepackage[most]{tcolorbox}
\tcbuselibrary{theorems,breakable,skins}
\usetikzlibrary{arrows.meta}

% Listas
\usepackage{enumitem}

% Cabeceras
\usepackage{fancyhdr}
\usepackage{emptypage}
\usepackage{titlesec}
\usepackage{titletoc}
\usepackage{makeidx}
\makeindex

% ==================================================
% COLORES
% ==================================================

\definecolor{maincolor}{HTML}{5B42F5}
\definecolor{darkcolor}{HTML}{29243D}
\definecolor{lightcolor}{HTML}{F4F2FF}
\definecolor{graycolor}{HTML}{666666}
\definecolor{linecolor}{HTML}{DCD9E8}

% Enlaces (hyperref se carga después de definir los colores)
\usepackage[
    colorlinks=true,
    linkcolor=maincolor,
    citecolor=maincolor,
    urlcolor=maincolor
]{hyperref}

% ==================================================
% FORMATO GENERAL
% ==================================================

\setlength{\parindent}{0pt}
\setlength{\parskip}{5pt}

\setcounter{secnumdepth}{3}
\setcounter{tocdepth}{2}

\setlist[itemize]{leftmargin=1.5em, itemsep=2pt, topsep=4pt}
\setlist[enumerate]{leftmargin=1.8em, itemsep=2pt, topsep=4pt}

% ==================================================
% CAPÍTULOS Y SECCIONES
% ==================================================

\linespread{1.06}

% (Los \par van dentro de macros porque titlesec no admite \par en sus argumentos)
\newcommand{\capfilete}{\textcolor{maincolor}{\rule{2.2cm}{1.2pt}}\par\vspace{1.2em}}
\newcommand{\capetiqueta}{%
  \textcolor{maincolor}{\scshape\large\chaptertitlename\ \thechapter}\par\vspace{0.5em}\capfilete}

% Capítulo numerado: etiqueta pequeña, filete corto y título grande
\titleformat{\chapter}[display]
{\normalfont\raggedright}
{\capetiqueta}
{0pt}
{\bfseries\fontsize{30}{34}\selectfont\color{darkcolor}}

% Capítulo sin numerar (Prefacio, Notación...)
\titleformat{name=\chapter,numberless}[display]
{\normalfont\raggedright}
{}{0pt}
{\capfilete\bfseries\fontsize{30}{34}\selectfont\color{darkcolor}}

\titlespacing*{\chapter}{0pt}{20pt}{34pt}

\titleformat{\section}
{\normalfont\Large\bfseries\color{darkcolor}}
{\textcolor{maincolor}{\thesection}}{0.8em}{}

\titleformat{\subsection}
{\normalfont\large\bfseries\color{darkcolor}}
{\textcolor{maincolor}{\thesubsection}}{0.8em}{}

% ==================================================
% ÍNDICE GENERAL
% ==================================================

\titlecontents{chapter}[2.5em]
{\addvspace{12pt}\bfseries\color{darkcolor}}
{\color{maincolor}\contentslabel{2.5em}}
{\hspace*{-2.5em}}
{\hfill\contentspage}

\titlecontents{section}[5.5em]
{}
{\contentslabel{3em}}
{\hspace*{-3em}}
{\titlerule*[0.6em]{.}\contentspage}

% ==================================================
% COMANDOS PERSONALIZADOS
% (añade aquí los tuyos; así los cambias en un solo sitio)
% ==================================================

\newcommand{\R}{\mathbb{R}}
\newcommand{\N}{\mathbb{N}}
\newcommand{\dd}{\mathrm{d}}
\newcommand{\e}{\mathrm{e}}
\newcommand{\dif}{\mathop{}\!\mathrm{d}}
\newcommand{\der}{\frac{\mathrm{d}}{\mathrm{d}t}}
\newcommand{\pd}{\partial}
\newcommand{\avg}[1]{\left\langle #1\right\rangle}

% Página preliminar suelta (cubierta, portada...) sin páginas en blanco intermedias
\newenvironment{pagina}{\clearpage\thispagestyle{empty}}{\clearpage}

% Separador ornamental dentro de un capítulo
\newcommand{\ornamento}{%
  \begin{center}\textcolor{maincolor}{\small$\bullet\ \ \bullet\ \ \bullet$}\end{center}}

% Epígrafe al inicio de un capítulo: \epigrafe{Texto}{Autor}
\newcommand{\epigrafe}[2]{%
  \begin{flushright}\begin{minipage}{0.7\textwidth}\raggedleft
  \small\itshape #1\par\vspace{3pt}
  {\upshape\textcolor{maincolor}{--- #2}}
  \end{minipage}\end{flushright}\vspace{1.5em}}

% ==================================================
% ENTORNOS (cajas con numeración automática)
% ==================================================
%
% Uso:
%   \begin{definicion}{Título}{etiqueta} ... \end{definicion}
%   -> se referencia con \ref{def:etiqueta}
%
% Todos comparten el mismo contador dentro de cada capítulo:
% Definición 1.1, Teorema 1.2, Ejemplo 1.3, ...

% Estilo común de las cajas con título
\tcbset{
    cajabase/.style={
        enhanced, breakable, arc=0pt,
        left=9pt, right=9pt, top=8pt, bottom=8pt,
        before skip=10pt, after skip=10pt,
        fonttitle=\bfseries, coltitle=darkcolor,
        attach boxed title to top left={yshift*=-\tcboxedtitleheight/2},
    }
}

% Definición
\newtcbtheorem[number within=chapter]{definicion}{Definición}{
    enhanced, breakable,
    colback=lightcolor, colframe=maincolor, colbacktitle=lightcolor,
    coltitle=darkcolor, fonttitle=\bfseries,
    boxrule=0pt, borderline west={2pt}{0pt}{maincolor},
    arc=0pt, left=9pt, right=9pt, top=8pt, bottom=8pt,
    before skip=10pt, after skip=10pt,
    detach title, before upper={\tcbtitle.\ \ },
}{def}

% Teorema
\newtcbtheorem[use counter from=definicion]{teorema}{Teorema}{
    enhanced, breakable,
    colback=white, colframe=linecolor, boxrule=0.5pt,
    borderline west={2pt}{0pt}{maincolor},
    arc=0pt, left=9pt, right=9pt, top=8pt, bottom=8pt,
    before skip=10pt, after skip=10pt,
    fonttitle=\bfseries, coltitle=darkcolor, colbacktitle=white,
    detach title, before upper={\tcbtitle.\ \ },
}{thm}

% Proposición
\newtcbtheorem[use counter from=definicion]{proposicion}{Proposición}{
    enhanced, breakable,
    colback=white, colframe=linecolor, boxrule=0.5pt,
    borderline west={2pt}{0pt}{maincolor},
    arc=0pt, left=9pt, right=9pt, top=8pt, bottom=8pt,
    before skip=10pt, after skip=10pt,
    fonttitle=\bfseries, coltitle=darkcolor, colbacktitle=white,
    detach title, before upper={\tcbtitle.\ \ },
}{prop}

% Lema
\newtcbtheorem[use counter from=definicion]{lema}{Lema}{
    enhanced, breakable,
    colback=white, colframe=linecolor, boxrule=0.5pt,
    borderline west={2pt}{0pt}{maincolor},
    arc=0pt, left=9pt, right=9pt, top=8pt, bottom=8pt,
    before skip=10pt, after skip=10pt,
    fonttitle=\bfseries, coltitle=darkcolor, colbacktitle=white,
    detach title, before upper={\tcbtitle.\ \ },
}{lema}

% Corolario
\newtcbtheorem[use counter from=definicion]{corolario}{Corolario}{
    enhanced, breakable,
    colback=white, colframe=linecolor, boxrule=0.5pt,
    borderline west={2pt}{0pt}{maincolor},
    arc=0pt, left=9pt, right=9pt, top=8pt, bottom=8pt,
    before skip=10pt, after skip=10pt,
    fonttitle=\bfseries, coltitle=darkcolor, colbacktitle=white,
    detach title, before upper={\tcbtitle.\ \ },
}{cor}

% Ejemplo
\newtcbtheorem[use counter from=definicion]{ejemplo}{Ejemplo}{
    enhanced, breakable,
    colback=white, colframe=linecolor, boxrule=0.5pt, arc=2pt,
    left=9pt, right=9pt, top=8pt, bottom=8pt,
    before skip=10pt, after skip=10pt,
    fonttitle=\bfseries, coltitle=darkcolor, colbacktitle=white,
    detach title, before upper={\tcbtitle.\ \ },
}{ej}

% Demostración (sin numeración; añade el cuadradito final sola)
\newtcolorbox{demostracion}[1][Demostración]{
    enhanced, breakable,
    colback=white, colframe=white, boxrule=0pt,
    borderline west={1pt}{0pt}{graycolor},
    arc=0pt, left=12pt, right=9pt, top=5pt, bottom=5pt,
    before skip=8pt, after skip=8pt,
    before upper={\textbf{#1.}\ },
    after upper={\hfill$\square$},
}

% Resultado importante (fondo de color)
\newtcolorbox{resultado}[1][]{
    enhanced, breakable,
    colback=maincolor, colframe=maincolor, coltext=white,
    boxrule=0pt, arc=2pt,
    left=10pt, right=10pt, top=9pt, bottom=9pt,
    before skip=12pt, after skip=12pt,
    #1
}

% Observación (sin numeración)
\newtcolorbox{observacion}[1][Observación]{
    enhanced, breakable,
    colback=gray!6, colframe=gray!6, boxrule=0pt,
    borderline west={2pt}{0pt}{graycolor},
    arc=0pt, left=9pt, right=9pt, top=7pt, bottom=7pt,
    before skip=8pt, after skip=8pt,
    before upper={\textbf{#1.}\ },
}

% ==================================================
% CABECERAS Y PIES
% ==================================================

\pagestyle{fancy}
\fancyhf{}

\renewcommand{\chaptermark}[1]{\markboth{#1}{}}
\renewcommand{\sectionmark}[1]{\markright{\thesection\ \ #1}}

\newcommand{\cabsep}{\hspace{1.2em}{\color{maincolor}\vrule width 0.5pt height 0.9em depth 0.25em}\hspace{1.2em}}

\fancyhead[LE]{\small\thepage\cabsep\scshape\leftmark}
\fancyhead[RO]{\small\scshape\rightmark\cabsep\thepage}

\renewcommand{\headrulewidth}{0.4pt}
\renewcommand{\headrule}{{\color{linecolor}%
  \hrule width\headwidth height\headrulewidth \vskip-\headrulewidth}}
\renewcommand{\footrulewidth}{0pt}

% Primera página de capítulo: solo número de página abajo
\fancypagestyle{plain}{%
  \fancyhf{}%
  \fancyfoot[C]{\small\thepage}%
  \renewcommand{\headrulewidth}{0pt}%
}

% Metadatos del PDF
\hypersetup{
  pdftitle={Mecánica de Distribuciones},
  pdfauthor={Mario Sultan Romero}
}


% --------------------------------------------------
% COMPILAR SOLO LO QUE ESTÁS ESCRIBIENDO
% --------------------------------------------------
% Descomenta la línea siguiente y deja solo los capítulos
% que quieras compilar. Es MUCHO más rápido en un libro grande.
% La numeración y las referencias del resto se conservan
% (hay que haber compilado el libro completo al menos una vez).
%
% \includeonly{capitulos/01-fundamentos}

\begin{document}

% --------------------------------------------------
% CUBIERTA Y PÁGINAS PRELIMINARES (sin numerar)
% --------------------------------------------------
\hypersetup{pageanchor=false}
\pagenumbering{alph}
% ==================================================
% CUBIERTA
% Fondo oscuro y una familia de gaussianas que se
% estrechan hasta convertirse en una delta de Dirac.
% ==================================================
\begin{pagina}
\thispagestyle{empty}
\begin{tikzpicture}[remember picture,overlay]

  % Fondo
  \fill[darkcolor] (current page.south west) rectangle (current page.north east);

  % Franja lateral fina
  \fill[maincolor] ([xshift=1.6cm]current page.south west)
        rectangle ([xshift=1.72cm]current page.north west |- current page.north west);

  % Gaussianas -> delta
  \begin{scope}[shift={([yshift=-4.2cm]current page.center)}]
    \draw[maincolor!50!white, opacity=0.35, line width=0.5pt]
      (-9,0) -- (9,0);
    \foreach \sig/\opa in {1.60/0.22, 1.20/0.30, 0.90/0.40, 0.65/0.52,
                           0.45/0.66, 0.30/0.82, 0.20/1.00}{
      \draw[color=maincolor!70!white, opacity=\opa, line width=0.9pt,
            domain=-9:9, samples=240, smooth, variable=\x]
        plot ({\x},{1.1/\sig*exp(-(\x*\x)/(2*\sig*\sig))});
    }
    % Delta de Dirac
    \draw[white, line width=1.1pt, -{Stealth[length=7pt,width=6pt]}] (0,0) -- (0,6.6);
  \end{scope}

  % Cabecera
  \node[anchor=north west, text=maincolor!35!white, inner sep=0pt]
    at ([xshift=2.8cm,yshift=-3cm]current page.north west)
    {\scshape\large Mecánica teórica};

  % Título
  \node[anchor=north west, text=white, align=left, inner sep=0pt, text width=15cm]
    at ([xshift=2.8cm,yshift=-4.6cm]current page.north west)
    {{\fontsize{44}{50}\selectfont\spaceskip=0pt\relax\bfseries Mecánica de\\[2pt] Distribuciones}\\[1.3cm]
     {\color{maincolor!45!white}\rule{3cm}{1.4pt}}\\[0.6cm]
     {\fontsize{14}{20}\selectfont\spaceskip=0pt\relax\color{white!85!maincolor}
      Fundamentos matemáticos\\ y formulación hamiltoniana}};

  % Autor
  \node[anchor=south west, text=white, inner sep=0pt]
    at ([xshift=2.8cm,yshift=2.4cm]current page.south west)
    {\scshape\Large Mario Sultan Romero};

\end{tikzpicture}
\end{pagina}
      % cubierta
% Anteportada (solo el título)
\begin{pagina}
\thispagestyle{empty}
\centering
\vspace*{0.3\textheight}
{\LARGE\bfseries\color{darkcolor} Mecánica de Distribuciones}
\vfill
\end{pagina}
   % anteportada (solo el título)
% Portada interior
\begin{pagina}
\thispagestyle{empty}
\centering

\vspace*{2.8cm}

{\large\color{maincolor}\scshape Mecánica teórica}

\vspace{1.6cm}

{\fontsize{34}{40}\selectfont\bfseries\color{darkcolor} Mecánica de\\[0.3cm] Distribuciones}

\vspace{0.9cm}

{\large\itshape Fundamentos matemáticos\\ y formulación hamiltoniana}

\vspace{1.6cm}

\textcolor{maincolor}{\rule{2.4cm}{1.2pt}}

\vfill

{\Large\scshape Mario Sultan Romero}

\vspace{1.2cm}

{\small 2026}

\vspace*{1cm}
\end{pagina}
       % portada interior
% Página de créditos (reverso de la portada)
\begin{pagina}
\thispagestyle{empty}
\vspace*{\fill}
\small\raggedright
\textbf{Mecánica de Distribuciones}\\
Fundamentos matemáticos y formulación hamiltoniana

\bigskip

\copyright\ 2026 Mario Sultan Romero.\\
Todos los derechos reservados.

\bigskip

Composición tipográfica realizada con \LaTeX.\\
Primera edición, 2026.\\
Versión del borrador: \today.

\bigskip

% \textbf{Contacto:} correo@ejemplo.com\\
% \textbf{Repositorio:} github.com/usuario/libro
\end{pagina}
      % créditos / copyright
% Dedicatoria
\begin{pagina}
\thispagestyle{empty}
\vspace*{0.3\textheight}
\begin{flushright}
\itshape
Para quien quiera\\
entender el porqué.
\end{flushright}
\vfill
\end{pagina}
   % dedicatoria

% --------------------------------------------------
% PÁGINAS INICIALES (numeración romana)
% --------------------------------------------------
\frontmatter
\hypersetup{pageanchor=true}
\chapter*{Prefacio}
\markboth{Prefacio}{Prefacio}
\addcontentsline{toc}{chapter}{Prefacio}

% Escribir aquí el prefacio.

\tableofcontents
\chapter*{Notación y convenciones}
\markboth{Notación y convenciones}{Notación y convenciones}
\addcontentsline{toc}{chapter}{Notación y convenciones}

A lo largo del libro se usa la siguiente notación.

\bigskip

\begin{tabular}{@{}p{3cm}p{9cm}@{}}
\toprule
\textbf{Símbolo} & \textbf{Significado} \\
\midrule
$\R,\ \N$ & Números reales y naturales \\
$\delta(x)$ & Delta de Dirac \\
$\rho(x,t)$ & Distribución espacial de partículas \\
$\avg{f}$ & Valor medio de $f$ \\
$\dif x$ & Elemento de integración \\
\bottomrule
\end{tabular}


\mainmatter

% --------------------------------------------------
% CAPÍTULOS  (uno por archivo; añade una línea por tema)
% --------------------------------------------------
\chapter{Fundamentos}
\label{cap:fundamentos}

% \epigrafe{Texto de la cita.}{Autor}

\section{Introducción}

% Introducción del capítulo.

\section{Distribución espacial}
\index{distribucion@distribución!espacial}

\begin{definicion}{Distribución de partículas}{distribucion-particulas}
Sea un sistema de $N$ partículas cuyas posiciones
vienen dadas por $x_i(t)\in\R^d$.
Definimos la distribución espacial normalizada como
\[
\rho(x,t)=\frac{1}{N}\sum_{i=1}^{N}\delta(x-x_i(t)).
\]
\end{definicion}

\begin{ejemplo}{Sistema de dos partículas}{dos-particulas}
Consideremos dos partículas situadas en $a$ y $b$.
La distribución correspondiente es
\[
\rho(x)=\frac12\delta(x-a)+\frac12\delta(x-b).
\]
Su normalización se verifica integrando sobre todo el espacio.
\end{ejemplo}

\begin{teorema}{Normalización}{normalizacion}
La distribución espacial de la Definición~\ref{def:distribucion-particulas}
satisface
\[
\int_{\R^d}\rho(x,t)\,\dif x=1.
\]
\end{teorema}

\begin{demostracion}
Sustituyendo la definición de la distribución,
\[
\int_{\R^d}\rho(x,t)\,\dif x
=\frac1N\sum_{i=1}^N\int_{\R^d}\delta(x-x_i(t))\,\dif x.
\]
Por la propiedad de normalización de la delta de Dirac,
cada integral vale la unidad. Por tanto,
\[
\int_{\R^d}\rho(x,t)\,\dif x=\frac1N\sum_{i=1}^N 1=1.
\]
\end{demostracion}

\begin{observacion}
La distribución no representa necesariamente una densidad ordinaria
respecto a la medida de Lebesgue. Debe interpretarse en sentido
distribucional.
\end{observacion}

% \include{capitulos/02-siguiente-tema}

% --------------------------------------------------
% APÉNDICES
% --------------------------------------------------
\appendix
\chapter{Resultados matemáticos auxiliares}
\label{ap:auxiliares}

% Escribir aquí los resultados auxiliares.


% --------------------------------------------------
% BIBLIOGRAFÍA E ÍNDICE ALFABÉTICO
% --------------------------------------------------
\backmatter

% \nocite{*} incluye todas las entradas de referencias.bib aunque no
% se citen. Bórralo cuando empieces a usar \cite{...} de verdad.
\nocite{*}
\phantomsection
\addcontentsline{toc}{chapter}{\bibname}
\bibliographystyle{plain}
\bibliography{referencias}

\phantomsection
\addcontentsline{toc}{chapter}{\indexname}
\printindex

\end{document}
